\documentclass[runningheads]{llncs}

\usepackage[T1]{fontenc}
\usepackage{url}
\usepackage{listings}

\lstset{
  basicstyle=\ttfamily\scriptsize,
  columns=fullflexible,
  keepspaces=true,
  breaklines=true,
  frame=single,
  xleftmargin=0.5em,
  xrightmargin=0.5em
}

\emergencystretch=2em

\title{An Executable Self-Justifying Axiom System in Proflog}
\titlerunning{Executable SJAS in Proflog}

\author{James Torre}
\authorrunning{J. Torre}
\institute{\email{jpt4@proton.me}}

\begin{document}
\maketitle

\begin{abstract}
Willard's self-justifying axiom systems (SJAS) are weak arithmetic
theories able to consistently assert selected formalizations of their own
consistency.  Their construction is intensional: the object language,
the coding of syntax, and the proof procedure jointly affect whether the
self-referential consistency tableau closes.  This system description reports
an implementation of an \(IS\#_D(\beta)\)-style SJAS in Proflog, Fitting's
tableau-based logic programming language.  The artifact supplies a finite SJAS
builder, binary U-grounding arithmetic, reflected axiom generation,
inspectable formula/system/proof codes, and proof predicates that consume
encoded tableau certificates.  Worked examples show ordinary Proflog execution
from frontend source to kernel proof traces, then show how reflected Proflog
clauses become part of the finite system whose consistency sentence is
generated and queried.
\keywords{self-justifying axiom systems \and logic programming \and semantic tableaux \and proof checking \and declarative programming}
\end{abstract}

\section{Research Objective}

Willard's self-verifying and self-justifying axiom systems mark a boundary in
the usual scope of G\"odel's second incompleteness theorem.  They do not make
Peano Arithmetic prove its own consistency.  They instead weaken the arithmetic
language so that a carefully formulated self-consistency sentence can be added
without inconsistency~\cite{willard2001selfverifying,willard2011selfjustifying}.

The programming-language question is whether this construction has an
executable counterpart.  If self-justification depends only on an extensional
set of theorems, a conventional prover for the same theorems might be enough.
If it depends on the intensional behavior of deduction, then a programming
language must preserve the relevant proof-tree structure, proof-code size, and
syntax-generation discipline.  This paper takes the second possibility
seriously.  It implements the ordinary-tableau \(IS\#_D(\beta)\) pattern in
Proflog and records where the implementation is structural, where it is merely
nominal, and where further correspondence work remains.

\section{Why Proflog}

\subsection{Self-Justification}

The relevant Willard systems use a weak arithmetic base usually described as
U-grounding.  The constants \(0\) and \(1\), successor-like growth, and addition
are available, but multiplication is not a total binary function.  Instead,
multiplication appears as a graph relation \(\mathit{mult}(x,y,z)\).  This
choice blocks enough arithmetic strength to avoid the standard second
incompleteness hypotheses while leaving enough coding capacity to discuss
formulae and proofs~\cite{willard2004addition,willard2013significance}.

For an \(IS\#_D(\beta)\) system, \(D\) is the deduction apparatus and
\(\beta\) is a finite list of proper axioms.  The self-consistency sentence
asserts that no \(D\)-proof derives the contradiction code from the finite
system code.  The exact point at which self-justification arises is therefore
not a model-theoretic slogan: it is a failure of a particular formalized proof
tree to close.  Tableau deduction matters because the generated tree branches
on the parent formula's structure, and the U-grounding language does not supply
the arithmetic strength needed to compress the pathological closure.

\subsection{Proflog}

Fitting's Proflog is a logic programming language whose operational semantics
is semantic tableaux with equality and a Procedure Call rule~\cite{fitting1994tableaux}.
To prove a query, the evaluator closes the tableau for the negated query.  To
refute a query, it closes the tableau for the query itself.  Procedure calls
open the compiled clause body, so ordinary program execution still produces
tableau proof terms rather than host-language truth values.

That makes Proflog a close executable target for the most linguistically
expressive SJAS variants based on semantic tableaux.  It is not enough that
Proflog and Willard's tableau method prove similar formulae.  A proof predicate
used inside an SJAS must preserve the relevant intensional measures: tree
branching, closure rules, proof constructors, and encoded proof size.  The
implementation below is therefore organized around visible descent from source
formulae to kernel proof certificates.

\section{Implementation Architecture}

The artifact repository is published at:

\begin{lstlisting}
https://github.com/distribution-artifacts/lopstr-ppdp26
\end{lstlisting}

It contains the paper sources, Proflog implementation, tests, and worked
examples cited below.

The implementation has three layers.  The AST and language layer builds
first-order formulae, terms, clauses, and signatures.  The frontend layer gives
users a compact Clojure-readable surface syntax and expands it into AST
constructors.  The kernel layer runs proof search and returns proof terms with
explicit constructors such as \texttt{split}, \texttt{witness},
\texttt{pos-call}, \texttt{neg-call}, \texttt{eq-step}, and closure rules.

\begin{lstlisting}
surface clauses/query
  -> frontend expansion
  -> AST clauses and formulae
  -> compiled Proflog program
  -> kernel tableau search
  -> proof term / answer / unresolved status
\end{lstlisting}

Proof profiles extend the kernel with domain-specific closure behavior while
keeping proof evidence visible.  The SJAS profile uses this hook for
U-grounding arithmetic, formula-code predicates, axiom membership, and encoded
proof checking.  The ordinary Proflog examples use the same kernel without the
SJAS profile.

\section{Worked Proflog Examples}

The first examples are Fitting's even/odd program and Nim program.  They are
small enough to inspect but exercise quantification, equality, recursive
procedure calls, and negative calls.

\begin{lstlisting}
(def p1-language
  (pf/language
    (constants zero)
    (functions (s 1))
    (relations (even 1) (odd 1))))

(def p1-program
  (pf/proflog p1-language
    (|- (even x)
        (or (= x zero)
            (exists [y]
              (and (= x (s y))
                   (odd y)))))
    (|- (odd x)
        (forall [y]
          (implies (even y)
                   (!= x y))))))
\end{lstlisting}

The catalog evaluator prints the kernel-level outcome and an abbreviated proof
trace.  The trace below is not a side calculation: it is the certificate shape
returned by the proof kernel.

\begin{lstlisting}
(fitting/evaluate-case :p1-odd-1-succeeds)
;; => {:outcome :succeeds,
;;     :proof-root neg-call,
;;     :proof-steps
;;     [neg-call witness conj eq-step par-bind
;;      pos-call split free-close ... refl-close]}
\end{lstlisting}

Fitting's Nim program is represented with the one-or-two-token move relation
inlined in the \texttt{win} clause:

\begin{lstlisting}
(def nim-program
  (pf/proflog p1-language
    (|- (win x)
        (exists [y]
          (and (or (= x (s y))
                   (= x (s (s y))))
               (not (win y)))))))
\end{lstlisting}

The recursive negative call leaves a visible proof signature.  Position 4 is a
winner; position 3 is refuted under the same kernel.

\begin{lstlisting}
(fitting/evaluate-case :p2-win-4-succeeds)
;; => {:outcome :succeeds,
;;     :proof-root neg-call,
;;     :proof-steps
;;     [neg-call once-univ split conj neq-close
;;      decompose args eq-bind pos-call witness ...]}

(fitting/evaluate-case :p2-win-3-fails)
;; => {:outcome :fails,
;;     :proof-root pos-call,
;;     :proof-steps
;;     [pos-call witness conj savefml split eq-step ...]}
\end{lstlisting}

Fitting warns that factoring the move test into an auxiliary relation changes
operational behavior.  The implementation preserves that warning: local
\texttt{move} queries are classified, but \texttt{win(1)} in the factored
program remains unresolved rather than being replaced by a host recurrence.

\begin{lstlisting}
(fitting/evaluate-case :factored-win-1-unresolved)
;; => {:outcome :unresolved,
;;     :classification :invalid-auxiliary-relation-factoring}
\end{lstlisting}

\section{Implementing \(IS\#_D(\beta)\) in Proflog}

\subsection{System Construction}

The SJAS builder accepts a proof profile, a finite \(\beta\), reflected user
clauses, and external user clauses.  Reflected clauses are compiled as ordinary
Proflog procedures and also become Group-2b axioms inside the generated finite
system.  External clauses are executable context only: they do not change the
system code or the generated self-consistency sentence.

\begin{lstlisting}
(def demo-system
  (sjas/system-source
    {:profile :willard-sjas-tableau0}
    (language
      (relations (demo 1)
                 (external-demo 1)))
    (beta
      (= 1 1))
    (reflected
      (|- (demo x)
          (= x 1)))
    (external
      (|- (external-demo x)
          (= x 0)))))

(frequencies (map :group (:axioms demo-system)))
;; => {:group-zero 2, :group-one 3, :group-two 1,
;;     :group-two-b 1, :group-three 1}
\end{lstlisting}

This distinction is important for self-reference.  A source clause is part of
the self-justified system only when it appears in the reflected finite basis
before Group-3 generation.  Merely calling an SJAS-flavored library from outside
does not enlarge the consistency claim.

\subsection{Codes}

The implementation exposes formula, system, and proof codes as first-order
terms.  The default compact representation writes byte strings with generated
\texttt{code-N} constructors.  The stronger U-grounding representation writes
the same byte strings as binary numerals over \(0\), \(1\), \texttt{dbl}, and
\texttt{add}, using a sentinel so trailing zero bytes remain injective.

\begin{lstlisting}
(take 12 (code/code-term-bytes (:system-code demo-system)))
;; => (31 32 2 5 25 1 0 1 25 1 0 1)

(:contradiction-code demo-system)
;; => (app code-1 (app dbl (app 1)))

(code/code-term-bytes (:contradiction-code demo-system))
;; => [2]
\end{lstlisting}

The finite symbol table is used to encode formula syntax, but the intended
invariant is structural up to signature isomorphism rather than nominal.
Regression tests rename reflected relation symbols, observe different byte
codes, and check that the proof certificate and proof-constructor size are
preserved.

\subsection{Proof Predicate}

The object-language predicate \texttt{tableau-proof(system,theorem,proof)}
checks an encoded certificate against an encoded finite system.  For the
implemented certificate classes it decodes the system code, reconstructs
axiom membership and reflected clauses from that code, decodes the theorem
code, decodes the proof-code tree, and checks the corresponding tableau,
equality, procedure-call, and axiom-citation constructors.  It does not decide
the predicate by consulting a host-side proof-target table.

\begin{lstlisting}
(sjas/query-succeeds
  demo-system
  (pf/q (demo 1))
  {:proof-limit 1 :fuel 160})
;; first proof steps:
;; [profiled willard-sjas-tableau0 conj neg-call
;;  profiled willard-sjas-arithmetic sjas-equal ...]

(let [beta-record (first (filter #(= :group-two (:group %))
                                 (:axioms demo-system)))
      cert (sjas/proof-certificate 'sjas-axiom)]
  (query/query-succeeds
    (:program demo-system)
    (sjas/tableau-proof (:system-code demo-system)
                        (:code beta-record)
                        cert)
    1
    200))
;; first proof steps:
;; [profiled willard-sjas-tableau0
;;  profiled willard-sjas-proof-check
;;  sjas-code-bytes willard-sjas-axiom-member
;;  sjas-system-beta-axiom ... sjas-axiom]
\end{lstlisting}

The Level-1 profile similarly exposes \texttt{subst-code/2} and
\texttt{subst-prf/4} for the fixed-point substitution path.  Open synthesis of
large public code numerals remains an operational boundary, so the focused
suite tests ground and partially bound decoding paths rather than claiming
unbounded code generation is cheap.

\subsection{In-Principle Arithmeticization}

The intended semantic object is not the fact that the host Proflog evaluator can
prove a formula.  It is the arithmetical relation
\(\mathsf{TabPrf}_\beta(s,t,p)\): the system code \(s\) decodes to a finite
\(IS\#_D(\beta)\) axiom basis \(\Gamma_s\), the theorem code \(t\) decodes to a
formula \(T\), and the proof code \(p\) decodes to a finite closed semantic
tableau for \(\Gamma_s \wedge \neg T\).  For the consistency sentence, \(t\)
is the contradiction code, so the Group-3 assertion says that no such \(p\)
exists for the system's own code.

This relation decomposes into bounded code readers, syntax checks, system-code
reconstruction, axiom-membership checking, proof-tree navigation, local tableau
rule checking, branch closure, substitution checking, and reflected-clause
expansion.  Each component is a relation over byte strings or U-grounding
numerals, expressible with the profile's arithmetic vocabulary.  Multiplication
need not be installed as a total function; fixed-radix decoding and other
arithmetic operations are used as graph relations on the concrete numerals that
appear in the codes.

Soundness is then the ordinary tableau soundness induction, but applied to the
decoded tree: every internal node is accepted only by a semantic-tableau rule
or by a reflected procedure-call macro equivalent to expanding a decoded finite
first-order clause from \(s\), and every leaf is accepted only by a branch
contradiction or by an arithmeticized profile relation.  Conversely, any finite
closed tableau using the admitted rules can be serialized as a proof code whose
local node evidence is accepted by the relation.  This is the in-principle
internalization required for the SJAS claim, even when evaluating the full
relation is impractical for large public numerals.

The proof-code layout is therefore required to be natural and inspectable, not
byte-for-byte identical to one historical presentation.  It must preserve the
tree and Willard's lower-bound size discipline for semantic tableaux; theorem
equivalence alone would be too weak, because a host prover could accept the
same formulas with shorter or differently shaped proof objects.

\section{Evaluation and Reproducibility}

The repository includes the commands needed to rebuild the paper and verify
the system examples:

\begin{lstlisting}
make paper
lein test-proflog-fast
lein test-proflog-extended
lein test-proflog-fitting-programs
lein run -m proflog.fitting-programs p1-odd-1-succeeds p2-win-4-succeeds p2-win-3-fails factored-win-1-unresolved
lein test :only proflog.willard-sjas-test/sjas-system-builder-generates-groups-and-reflected-boundary
lein test :only proflog.willard-sjas-test/sjas-u-grounding-tableau-proof-checks-numeral-system-theorem-and-proof-codes
\end{lstlisting}

In the artifact repository, the focused fast suite passed 159 tests / 594
assertions in 146.92 seconds, the extended suite passed 68 tests / 203
assertions in 267.79 seconds, and the Fitting catalog gate passed 6 tests / 81
assertions in 71.02 seconds after syncing with the parent repository's current
result surface.  The corresponding parent Fitting run passed in 72.23 seconds.

The SJAS namespace and the full Fitting catalog are intentionally expensive.
The artifact therefore provides progress-visible SJAS aliases and direct
catalog-case execution so a reader can see progress one semantic test at a time
before running a whole namespace.  This matters for claims about proof
predicates: a stalled full run gives little evidence, while focused tests
identify whether code decoding, axiom membership, proof-code decoding,
U-grounding arithmetic, or reflected procedure calls are responsible.

\section{Limitations and Future Work}

The implementation is an executable system description, not a completed
formal correspondence proof between Proflog and every Willard tableau
presentation.  A theorem-level equivalence would be too weak for SJAS work: a
host prover could prove the same theorems with shorter or differently shaped
certificates and thereby change the self-referential consistency behavior.  The
relevant future theorem must preserve the tableau-structural invariants on
which self-justification depends, while proving irrelevant differences harmless
up to isomorphism.

The current proof predicate covers the certificate constructors exercised by
the generated SJAS examples and regressions.  Further work includes broader
proof-code synthesis, stronger admissibility checks for reflected Group-2b
extensions, performance improvements in recursive proof search, and a formal
definition of the Proflog kernel suitable for comparing
\(IS\#_{D_{\mathrm{Proflog}}}(\beta)\) with Willard's semantic-tableau
\(IS\#_D(\beta)\).  A more speculative direction is computationalizing
unclosable tableau branches as nonterminating or complexity-barrier programs.

\section{Related Work}

The mathematical basis is Willard's work on self-verifying and
self-justifying axiom systems~\cite{willard2001selfverifying,willard2011selfjustifying}
and his analysis of semantic tableaux near Robinson arithmetic~\cite{willard2004addition,willard2013significance}.
The programming-language basis is Fitting's Proflog~\cite{fitting1994tableaux}.
The implementation uses Clojure and \texttt{core.logic}~\cite{corelogic}, but
the kernel records its own proof terms rather than treating host unification as
the public proof object.

\section{Conclusion}

Proflog gives a practical executable setting for studying Willard-style
self-justification because its ordinary programming semantics is already
tableau proof search.  The artifact described here implements a finite
\(IS\#_D(\beta)\)-style builder, U-grounding arithmetic, reflected Proflog
extensions, inspectable syntax/proof codes, and encoded proof predicates.  The
result is useful both as a working logic-programming system and as a concrete
testbed for the harder question: which intensional features of a deductive
apparatus must be preserved for self-justification to remain meaningful?

\paragraph{Acknowledgements.}
This work was supported by the FUTO Fellowship, Summer 2025; Bloominglabs,
Inc.; and Seth Frey, Department of Communication, UC Davis.

\paragraph{AI disclosure.}
Project material was created in part using Codex GPT-5.4 and GPT-5.5
(OpenAI), and Claude Opus 4.6 and 4.7 (Anthropic).

\bibliographystyle{splncs04}
\bibliography{references}

\end{document}
